\documentclass{extarticle}

\usepackage{lineno}
\usepackage[T2A]{fontenc}
\usepackage[utf8]{inputenc}
\usepackage[russian]{babel}
\usepackage{amsmath, amsfonts, amssymb}
\usepackage{float}
\usepackage{wrapfig}
\usepackage{graphicx}
\usepackage{enumitem}
\usepackage{csquotes}
\usepackage{hyperref}
\usepackage[margin=2cm]{geometry}
\usepackage{setspace}
\usepackage{indentfirst}
\usepackage{subcaption}
\usepackage[style=gost-numeric, backend=biber]{biblatex}

\usepackage[ruled,vlined,linesnumbered]{algorithm2e}

\addbibresource{references.bib}

\begin{document}
    \begin{titlepage}
        \newpage
        \begin{center}
        
            \bfseries ФЕДЕРАЛЬНОЕ ГОСУДАРСТВЕННОЕ БЮДЖЕТНОЕ ОБРАЗОВАТЕЛЬНОЕ УЧРЕЖДЕНИЕ ВЫСШЕГО ОБРАЗОВАНИЯ <<МОСКОВСКИЙ ГОСУДАРСТВЕННЫЙ УНИВЕРСИТЕТ имени М.\@В.\@ ЛОМОНОСОВА>>
            
            \vspace{1cm}
            ФИЗИЧЕСКИЙ ФАКУЛЬТЕТ
            \vspace{1cm}
        
            КАФЕДРА ФИЗИКИ ЭЛЕМЕНТАРНЫХ ЧАСТИЦ
        
            \vspace{6cm}
            ОТЧЕТ ПО НАУЧНО-ИССЛЕДОВАТЕЛЬСКОЙ РАБОТЕ 
            
            НАПРАВЛЕНИЕ 03.04.02 \"ФИЗИКА\"
            
            СЕМЕСТР 09
        \end{center}

        \vspace{4cm}

        \begin{flushright}
            Выполнил:
            
            студент 5 курса 509 группы
            
            Балинов Темир Максимович
        \end{flushright}
        
        \vspace{0.5cm}

        \begin{flushright}
            Научный руководитель:
            
            к.ф.-м.н.
            
            Апарин Алексей Андреевич

            Научный консультант:
            
            Б.В. Лыонг
        \end{flushright}

        \vspace{\fill}

        \begin{flushleft}
            Отчет принят: 

            \vspace{0.5cm}

            Руководитель НИР:
        \end{flushleft}
        
        \begin{center}
            Дубна 2026
        \end{center}

    \end{titlepage}
\tableofcontents
\newpage
\linenumbers{}
\section{Введение}
    Изучение свойств сильно взаимодействующей материи в условиях экстремальных плотностей и температур 
    является одной из центральных задач современной физики высоких энергий. 
    Особый интерес представляет область низких и средних энергий столкновений тяжёлых ионов 
    ($\sqrt{s_{NN}} \sim 2$--$10$~ГэВ), где, согласно теоретическим предсказаниям, 
    может проявляться фазовый переход первого рода и существовать критическая точка квантовой хромодинамики (КХД).

    Фемтоскопия идентичных частиц~\cite{Lisa:2005dd} предоставляет уникальный метод 
    исследования пространственно-временной структуры источника частиц на момент кинетического вымораживания. 
    Измеряя корреляции идентичных пионов, можно извлечь размеры области испускания 
    ($R_{\text{out}}$, $R_{\text{side}}$, $R_{\text{long}}$) и параметр хаотичности $\lambda$, 
    чувствительные к уравнению состояния и динамике расширения системы.
    
    \subsection{Цели и задачи работы}

        \textbf{Цель работы:} исследование пространственно-временных характеристик источника пионов 
        в столкновениях $^{197}\text{Au} + ^{197}\text{Au}$ при $\sqrt{s_{NN}} = 3$~ГэВ на основе данных модели SMASH 
        методом кореляционной фемтоскопии.

        \textbf{Для достижения цели поставлены следующие задачи:}
        \begin{enumerate}[leftmargin=*]
            \item Разработка и верификация конвертера данных из формата OSCAR2013 
                в формат ROOT (McDst) с корректным извлечением координат кинетического вымораживания.
            
            \item Исследование зависимостей извлечённых радиусов и параметра хаотичности 
                $\lambda$ от быстроты в интервале $y \in [-1; 1]$.
            
            \item Построение отношения проекций корреляционных функций 
                для $\pi^+$ и $\pi^-$ с целью поиска признаков зарядовой асимметрии.
        \end{enumerate}
\newpage
    
\section{Теоретические основы фемтоскопии}
    Основным инструментом исследования пространственно-временных характеристик 
    области испускания частиц является двухчастичная корреляционная функция. 
    В основе метода лежит квантовов статистический формализм Копылова-Подгорецкого, 
    описывающий корреляции идентичных бозонов, испускаемых хаотическим источником.

    Амплитуда вероятности обнаружить частицу с 4-координатой \textbf{x} от источника А 
    в точке $\mathbf{x_{A}}$ зависит только от разности $(\mathbf{x} - \mathbf{x_{A}})$:

    \begin{equation}
        \langle x|\psi_{A}\rangle =
            {(2\pi)}^{-4}
            \int 
                d^4\kappa u_{A}(\kappa) 
                \exp(i\kappa(\mathbf{x} - \mathbf{x_{A}})) 
    \end{equation}

    Где:
    \begin{align*}
        &\kappa = (\kappa_0, \mathbf{\kappa}) 
            - \text{4-импульс в пространстве фурье} \\
        &u_{A}(\kappa) 
            - \text{амплитуда вероятности того, что источник А испустит частицу с 4-импульсом }\kappa \\
        &\exp(i\kappa(\mathbf{x} - \mathbf{x_{A}})) 
            - \text{плоская волна, распространяющаяся из точки } x_A
    \end{align*}

    Пользуясь полнотой базиса координатных состояний перейдем 
    к 4-импульсному представлению через $\langle p|x\rangle = \exp(-ipx)$, получаем:
    \begin{equation}
        \langle p | \psi_{A}\rangle = 
            u_{A}(p) \exp(-ipx_{A})
    \end{equation}

    Для двух идентичных бозонов со спином 0 
    полная волновая функция должна быть симметричной. 
    Тогда амплитуду можно записат следующим образом:

    \begin{equation}
        T^{sym}_{AB} (p_{1}, p_{2}) = 
        \frac{
            \langle p_{1} | \psi_{A}\rangle
            \langle p_{2} | \psi_{B}\rangle + 
            \langle p_{2} | \psi_{A}\rangle 
            \langle p_{1} | \psi_{B}\rangle
        }{ 
            \sqrt{2}
        }
    \end{equation}

    Тогда, поскольку корреляционная функция --- 
    отношение двучастичного инклюзивного спектра 
    к произведению одночастичных инклюзивных спектров.

    \begin{align}
        C(p_1, p_2) &= 
            \frac{
                P_2(p_1, p_2)
            }{
                P_1(p_1)P_1(p_2)
            }
    \end{align}
    Где,
    \begin{align}
        P_2(p_1, p_2) = 
        |T^{sym}_{AB}|^2 &= 
        \frac{1}{2} 
        \bigg[
            |\langle p_{1} | \psi_{A} \rangle|^2 
            |\langle p_{2} | \psi_{B}\rangle|^2 + 
            |\langle p_{2} | \psi_{A} \rangle|^2 
            |\langle p_{1} | \psi_{B}\rangle|^2 + \nonumber \\
            & \quad + \mathfrak{R}
            \left(
                \langle p_{1} | \psi_{A}\rangle 
                \langle p_{2} | \psi_{B}\rangle 
                \langle p_{2} | \psi_{A}\rangle^{*} 
                \langle p_{1} | \psi_{B}\rangle^{*}
            \right)
        \bigg] \nonumber \\
        &= |u_{A}(p_1)|^2 |u_{B}(p_2)|^2 + 
        |u_{A}(p_2)|^2 |u_{B}(p_1)|^2 + \nonumber \\
        & \quad + 
        \mathfrak{R}(
            u_{A}(P_1)
            u_{B}(p_2)
            u^{*}_{A}(p_2)
            u^{*}_{B}(p_1) 
            \exp(-i(p_1 - p_2)(x_{A} - x_{B}))
        )\\
        P_1(p) &= 
            \sum_{A} | \langle p|\psi_{A}\rangle|^2 = 
            \sum_{A} |u_{A}(p)|^2 
    \end{align}

    Обозначим:
    \begin{align*}
        &q = p_{1} - p_{2} 
            - \text{относительный импульс} \\
        &\Delta x = x_{A} 
            - x_{B} - \text{разность координат источника}
    \end{align*}
    Также усредним по всем парам источников $A$ и $B$:

    \begin{equation}
        C (p_1, p_2) = 1 + \frac{
            \mathfrak{R}
            \sum_{AB} 
                u_A(p_1)
                u_B(p_2)
                u^{*}_A(p2)
                u^{*}_B(p_1) 
                \exp(-iq\Delta x)
        }{
            \sum_{AB} 
                |u_A(p_1)u_B(p_2)|^2
        }
    \end{equation}

    Допустим приближение гладкости:

    \begin{align*}
        &u_A(p_1) \approx u_A(p2) 
            - \text{амплитуды медленно меняются с иммпульсом} \\
        & \text{Тогда:} 
            \quad u_A(p_1) u^{*}_{A}(p_2) \approx |u_A(p)|^2
    \end{align*}

    А корреляционная функция упростится:

    \begin{equation}
        C(p_1, p_2) \approx 1 + \langle \cos(q\Delta x)\rangle
    \end{equation}

    Где усреднение проводится по всем парам источника.

    Эта формула содержит в себе информацию о пространственном распределении источника ($\Delta x$) 
\newpage

\section{Модель SMASH}
    Для моделирования динамики столкновений тяжёлых ионов 
    в данной работе используется программный комплекс \textbf{\texttt{SMASH}} 
    (Simulating Many Accelerated Strongly-interacting Hadrons)~\cite{Weil:2016zrk}. 
    \texttt{SMASH} представляет собой микроскопический транспортный подход, 
    предназначенный для описания эволюции адронной материи в условиях неравновесных процессов.

    \subsection{Физическая основа}
        В основе модели лежит решение релятивистского кинетического уравнения Больцмана 
        для функции распределения частиц $f_i(x, p)$ в фазовом пространстве~\cite{Weil:2016zrk}:
        \begin{equation}
            p^\mu \partial_\mu f_i(x, p) = C_i[f],
        \end{equation}
        где левая часть уравнения описывает свободное релятивистское движение частиц, 
        а правая часть $C_i[f]$ представляет собой интеграл столкновений, 
        учитывающий взаимодействия между частицами.

        Транспортное уравнение решается методом Монте-Карло 
        с использованием подхода тестовых частиц. 
        Эволюция системы моделируется как последовательность 
        свободных пробегов и бинарных столкновений~\cite{Weil:2016zrk}:
        \begin{itemize}
            \item \textbf{Свободное движение:} 
                частицы перемещаются по классическим траекториям 
                в течение малого шага по времени $\Delta t$;
            \item \textbf{Столкновения:} 
                при сближении частиц на расстояние, 
                определяемое сечением взаимодействия $\sigma_{tot}$, 
                происходит рассеяние или рождение новых частиц;
            \item \textbf{Распады:} 
                нестабильные резонансы распадаются согласно 
                экспоненциальному закону с учётом релятивистского замедления времени.
        \end{itemize}

    \subsection{Спектр частиц и взаимодействия}
        Модель оперирует адронными степенями свободы и включает 
        все хорошо установленные адроны с массой до $\sim 2$~ГэВ~\cite{Weil:2016zrk}:
        \begin{itemize}
            \item \textbf{Мезоны:} 
                $\pi$, $K$, $\eta$, $\rho$, $\omega$, $\phi$ и их возбуждённые состояния;
            \item \textbf{Барионы:} 
                нуклоны ($N$), $\Delta$-резонансы, гипероны ($\Lambda$, $\Sigma$, $\Xi$, $\Omega$) и их резонансы.
        \end{itemize}

        Сечения взаимодействий и вероятности каналов распада взяты из экспериментальных данных PDG 
        и параметризаций, используемых в других транспортных моделях (в частности, UrQMD)~\cite{Weil:2016zrk}. 
        При низких энергиях доминируют процессы возбуждения и распада резонансов, 
        тогда как при высоких энергиях включаются механизмы струнной фрагментации.

    \subsection{Области применения}
        Согласно~\cite{Weil:2016zrk}, модель \texttt{SMASH} 
        прошла всестороннюю проверку и рекомендуется для использования в следующих режимах:
        \begin{itemize}
            \item Столкновения тяжёлых ионов при кинетических энергиях пучка 0.4--2.0~AГэВ;
            \item Описание адронной стадии в гибридных моделях при высоких энергиях ($\sqrt{s_{NN}} > 7.7$~ГэВ);
            \item Расчёты свойств бесконечной ядерной материи в равновесии.
        \end{itemize}

        Модель обеспечивает сохранение энергии, импульса и квантовых чисел в каждом акте взаимодействия, 
        что делает её пригодной для изучения транспортных свойств адронной материи и сравнения с 
        экспериментальными данными установок SPS, RHIC, NICA и FAIR~\cite{Weil:2016zrk}.


    \subsection{Применение модели \texttt{SMASH} в данной работе}
        В рамках настоящей работы проведена серия расчётов столкновений 
        ядер золота $^{197}\text{Au} + ^{197}\text{Au}$ 
        с использованием генератора событий SMASH версии 3.1~\cite{SMASH:Code}. 
        Ниже приведены детали конфигурации и параметры моделирования.

        \subsubsection{Параметры столкновения}
            Основные характеристики моделируемой системы 
            представлены в таблице~\ref{tab:smash_params}.

            \begin{table}[H]
                \centering
                \caption{
                    Параметры конфигурации SMASH 
                    для расчётов Au+Au при $\sqrt{s_{NN}} = 3$~ГэВ
                }\label{tab:smash_params}
                \begin{tabular}{|l|l|}
                    \hline
                    \textbf{Параметр} & \textbf{Значение} \\ \hline
                    Система & $^{197}\text{Au} + ^{197}\text{Au}$ \\ \hline
                    Энергия ($\sqrt{s_{NN}}$) & 3~ГэВ \\ \hline
                    Состав ядра & 79 протонов + 118 нейтронов \\ \hline
                    Прицельный параметр ($b$) & 0.0--16.0~фм \\ \hline
                    Система отсчёта & Центр масс \\ \hline
                    Ферми-движение & frozen \\ \hline
                \end{tabular}
            \end{table}

        \subsubsection{Параметры эволюции системы}
            Временная эволюция системы моделировалась 
            со следующими параметрами~\cite{Weil:2016zrk}:
            \begin{itemize}
                \item \textbf{Шаг по времени:} 
                    $\Delta t = 0.1$~фм/$c$ (фиксированный режим);
                \item \textbf{Время эволюции:} 
                    $t_{\text{end}} = 200.0$~фм/$c$;
                \item \textbf{Принудительный распад:} 
                    включён (\texttt{Force\_Decays\_At\_End = true}).
            \end{itemize}

            Опция принудительного распада гарантирует, что все нестабильные резонансы, 
            оставшиеся к моменту окончания симуляции, распадаются на стабильные частицы. 
            Это обеспечивает корректный учёт вкладов от резонансов 
            в конечные спектры и сохранение законов сохранения~\cite{Weil:2016zrk}.

\section{Практические результаты работы}
    \subsection{Проверка корректности извлечение информации о кинетическом вымораживании}
        %%%%%%%%%%%%%%%
        %%%%%%%%%%%%%%% дополнить текст, типо , чтобы сравнивать нужно проанализировать одинаковость извлеченных импульсов и энергии с готовым файлом
        %%%%%%%%%%%%%%%

        \subsubsection{Описание алгоритма}
            Для анализа результатов требовалось преобразовать данные 
            из стандартного для \texttt{SMASH} формата (\texttt{OSCAR2013}) 
            в формат \texttt{ROOT} (\texttt{McDst}) с необходимой для нас обработкой, 
            включающей выставление координат кинетического вымораживания.

            Стандартный формат \texttt{OSCAR2013} для \texttt{SMASH} имеет следующую структуру:

            \begin{verbatim}
                # ... (служебная информация: версия генератора, структура файла)
                # event <ev_num> in 394
                ... (список начальных частиц до первого столкновения)
                # interaction in <in_part> out <out_part> ...
                ... (частицы, участвующие во взаимодействии)
                # event <ev_num> out <n_part>
                ... (список частиц на момент freeze-out, t = 200 fm/c)
                # event <ev_num> end 0 impact <b> ...
                ... (метаданные: прицельный параметр b, тип рассеяния)
            \end{verbatim}

            Для его преобразования использовался следующий алгоритм:

            \begin{itemize}
                \item \textbf{Построчный парсинг:} 
                    каждая строка файла проверяется на наличие символа \texttt{'\#'}. 
                    Строки, начинающиеся с этого символа, 
                    интерпретируются как служебные заголовки блоков, 
                    остальные строки содержат данные о частицах.

                \item \textbf{Обработка блока взаимодействий:} 
                    при обнаружении ключевых слов \texttt{\# interaction} 
                    начинается чтение частиц, участвовавших в столкновениях. 
                    Все уникальные частицы (по идентификатору \texttt{ID}) 
                    сохраняются в буфер \texttt{interactions}. 
                    При повторении \texttt{ID} перезаписывается информация 
                    о последней встреченной частице.

                \item \textbf{Фильтрация событий:}
                    при чтении заголовка \texttt{\# event \ldots out} 
                    анализируется количество частиц \texttt{nPart}. 
                    Если \texttt{nPart = 394} (сумма нуклонов в двух ядрах $^{197}\text{Au}$), 
                    событие классифицируется как упругое и исключается из дальнейшего анализа.

                \item \textbf{Классификация частиц:} 
                    данные о частицах распределяются по буферам в зависимости от контекста:
                    \begin{itemize}
                        \item 
                            если чтение происходит после \texttt{interaction} 
                            --- частица записывается в буфер \texttt{interactions};
                        \item 
                            если чтение происходит после \texttt{event \ldots out} 
                            --- частица записывается в буфер \texttt{out}.
                    \end{itemize}

                \item \textbf{Финализация события:} 
                    при чтении заголовка \texttt{\# event \ldots end} 
                    производится окончательная обработка частиц из буфера \texttt{out}:
                    \begin{itemize}
                        \item \textbf{Спектаторы:} 
                            частицы, присутствующие в конечном состоянии, 
                            но отсутствующие в буфере \texttt{interactions} 
                            (при условии $\text{ID} < 394$), помечаются как спектаторы. 
                            Их координаты и время устанавливаются в \texttt{-999} 
                            для исключения из фемтоскопического анализа;
                        \item \textbf{Участники:} 
                            для частиц, участвовавших во взаимодействиях, сохраняются 
                            координаты кинетического вымораживания ($t \approx 200$~фм/$c$). 
                            Эти данные записываются в итоговый файл формата McDst.
                    \end{itemize}
            \end{itemize}

            Подробный алгоритм описанный на языке псевдокода приведен ниже.

            \begin{algorithm}[H]
                \small
                \caption{Идея алгоритма конвертации OSCAR $\to$ McDst}\label{alg:pseudocode}

                \KwIn{Файл SMASH (OSCAR format)}
                \KwOut{ROOT-файл с деревом McDst}

                \textbf{Константы:} 
                    $N_{\text{init}} = 394$, $t_{\text{freeze}} = 200$~фм/$c$\;
                \textbf{Инициализация:} 
                    $\text{interactions} \leftarrow \emptyset$, 
                    $\text{out} \leftarrow \emptyset$, 
                    $\text{isElastic} \leftarrow \text{false}$\;
                    \BlankLine
                \ForEach{строка $L$ во входном файле}{
                    \If{$L$ начинается с \texttt{'\#'}}{
                        \If{$L$ содержит \texttt{\"interaction\"}}{
                            \While{чтение строк с частицами}{
                                извлечь частицу $p$\;
                                $\text{interactions}[p.\text{ID}] \leftarrow p$\;
                            }
                        }
                        \ElseIf{$L$ содержит \texttt{\"event \ldots out\"}}{
                            распарсить $nPart$\;
                            \If{$nPart == N_{\text{init}}$}{
                                $\text{isElastic} \leftarrow \text{true}$\;
                            }
                        }
                        \ElseIf{$L$ содержит \texttt{\"event \ldots end\"}}{
                            распарсить $\text{impactParameter}$\;
                            \If{$\text{isElastic} == \text{true}$}{
                                $\text{impactParameter} \leftarrow -1$\;
                            }
                            \If{$\text{isElastic} == \text{false}$}{
                                \ForEach{частица $p \in \text{out}$}{
                                    записать $p$ в McDst-дерево\;
                                }
                            }
                            $\text{interactions}.\text{clear}()$\;
                            $\text{out}.\text{clear}()$\;
                            $\text{isElastic} \leftarrow \text{false}$\;
                        }
                        \textbf{перейти к следующей строке}\;
                    }
                
                    \If{$\text{isElastic} == \text{true}$}{
                        \textbf{перейти к следующей строке}\;
                    }
                
                    извлечь параметры частицы $p$ из $L$\;
                
                    \eIf{$p.\text{ID} \notin \text{interactions}$ \textbf{и} $p.\text{ID} < N_{\text{init}}$}{
                        $p.\text{type} \leftarrow \text{SPECTATOR}$\;
                        $\text{out}[p.\text{ID}] \leftarrow p$\;
                    }{
                        \If{$p.\text{ID} \in \text{interactions}$}{
                            $p_{\text{final}} \leftarrow \text{interactions}[p.\text{ID}]$\;
                            $p_{\text{final}}.\text{coords} \leftarrow p.\text{coords}$\;
                            $\text{out}[p.\text{ID}] \leftarrow p_{\text{final}}$\;
                        }
                        \Else{
                            $\text{out}[p.\text{ID}] \leftarrow p$\;
                        }
                    }
                }

                \textbf{Финализация:} записать мета-информацию, сохранить файл\;
            \end{algorithm}

        \subsubsection{Проверка корректности конвертера}

            Для валидации алгоритма конвертации было сгенерировано тестовое событие 
            с сидом \texttt{2884011568699097786} в двух форматах: \texttt{Root} и \texttt{Oscar2013}. 
            Файл в формате \texttt{Root} использовался как эталонный результат генератора \texttt{SMASH}, 
            тогда как файл \texttt{Oscar2013} подвергался обработке описанным выше конвертером.

            Сравнение проводилось для распределений компонент импульса 
            $p_{x}$, $p_{y}$, $p_{z}$ и энергии $E$. 
            Выбор этих величин обусловлен тем, что конвертер модифицирует 
            исключительно пространственные координаты частиц 
            (заменяя их на координаты кинетического вымораживания), 
            не затрагивая их импульсно-энергетические характеристики. 
            Следовательно, соответствующие распределения должны совпадать.

            \begin{figure}[H]
                \centering
                \includegraphics[width=0.75\linewidth]{images/plots_mcdst.png}
                \caption{
                    Распределения $p_{x}$, $p_{y}$, $p_{z}$ и $E$, 
                    полученные после обработки файла \texttt{Oscar2013} конвертером (формат \texttt{McDst}).
                }\label{fig:mcdst_distributions}
            \end{figure}

            \begin{figure}[H]
                \centering
                \includegraphics[width=0.75\linewidth]{images/plots_smash.png}
                \caption{
                    Эталонные распределения 
                    $p_{x}$, $p_{y}$, $p_{z}$ и $E$, 
                    полученные непосредственно из генератора SMASH (формат \texttt{Root}).
                    }\label{fig:smash_distributions}
            \end{figure}

            \begin{figure}[H]
                \centering
                \includegraphics[width=0.75\linewidth]{images/plots_ratio.png}
                \caption{
                    Отношение распределений \texttt{McDst}/\texttt{SMASH}.
                    Значения, равные  единице, свидетельствуют о корректности конвертации.
                }\label{fig:ratio_distributions}
            \end{figure}

            Визуальный анализ гистограмм 
            (рис.~\ref{fig:mcdst_distributions}--\ref{fig:smash_distributions}) 
            показывает полное совпадение распределений. 
            Отношение гистограмм (рис.~\ref{fig:ratio_distributions}) 
            демонстрирует значения, равные единице по всему диапазону, 
            что подтверждает корректность преобразования данных из формата 
            \texttt{Oscar2013} в формат \texttt{Root} (\texttt{McDst}) 
            без искажения импульсно-энергетических характеристик частиц.

    \subsection{Составление и анализ зависимотси радиусов и силы корреляции от быстроты}
        Для анализа выполняются следующие условия:
        \begin{itemize}
            \item Частицы разделяются по знаку заряда: 
                $\pi^+$ и $\pi^-$;
            \item Данные разделяются на классы центральности: 
                $0\%$--$10\%$, $10\%$--$30\%$;
            \item Данные разделяются на классы по быстроте на диапозоне 
                $y \in [-1;1]$ с шагом 0.2;
            \item Ограничения на пары частиц kt:
                $\in$ [0.15 \- 0.6];
            \item Ограничения на одночастичные треки: 
                $p_T \in [0.15, 0.6]$ а также 
                $p \in [0.15, 0.6]$;
        \end{itemize}

        А также в качестве аппроксимирующей выбирается следующая функция: 
        \begin{equation*}
            CF(q_{out}, q_{side}, q_{long}) = 
                1 + \lambda \exp(
                    \frac{
                        -(
                            R_{out}^2q_{out}^2 +
                            R_{side}^2q_{side}^2 +
                            R_{long}^2q_{long}^2 +
                            2R_{out-long}q_{out}q_{long}
                        )
                    }{{(hc)}^{2}}
                )
        \end{equation*}

        \subsubsection{Аппроксимация гистограмм}
            Аппроксимация производится в диапозоне $q \in [-0.2, 0.2]$. 
            Ниже, приведены 1-мерные проекции 3-мерной 
            кореляционной функции в области центральных быстрот:

            \begin{figure}[H]
                \centering
                \includegraphics[width=0.8\linewidth]{cfs_pos_0-10_0.00_0.20.pdf}
                \caption{
                    Проекции кореляционных функций с наложенным фитом для $\pi^+$
                }
            \end{figure}

            \begin{figure}[H]
                \centering
                \includegraphics[width=0.8\linewidth]{cfs_neg_0-10_0.00_0.20.pdf}
                \caption{
                    Проекции кореляционных функций с наложенным фитом для $\pi^-$
                }
            \end{figure}

        \subsubsection{Анализ зависимости радиусов и силы корреляции}

            Получив параметры из апроксимации кореляционных функций, 
            были построены следующие графики зависимости радиусов и силы корреляции от быстроты.

            \begin{figure}[H]
                \centering
                \begin{subfigure}[b]{\textwidth}
                    \includegraphics[width=\textwidth]{images/c_radii_pos.pdf}
                \end{subfigure}
                \hfill
                \begin{subfigure}[b]{0.7\textwidth}
                    \includegraphics[width=\textwidth]{images/c_cross_and_lambda_graphs_pos.pdf}
                \end{subfigure}
                \caption{
                    Зависимость радиусов и силы корреляциии от диапазона y ($\pi^{+}$)
                }
            \end{figure}

            \begin{figure}[H]
                \centering
                \begin{subfigure}[b]{\textwidth}
                    \includegraphics[width=\textwidth]{images/c_radii_neg.pdf}
                \end{subfigure}
                \hfill
                \begin{subfigure}[b]{0.7\textwidth}
                    \includegraphics[width=\textwidth]{images/c_cross_and_lambda_graphs_neg.pdf}
                \end{subfigure}
                \caption{
                    Зависимость радиусов и силы корреляциии от диапазона y ($\pi^{-}$)
                }
            \end{figure}

            Визуальный анализ приведенных графиков показывает, 
            что для области центральных быстрот радиус источника 
            будет является наименьшим возможным, а также, 
            что при возрастании модуля быстроты будет расти и радиус источника. 
            Можно заметить, что похожее поведение имеет и график для $\lambda$.

            Также, анализ показывает, рост недиагонального радиуса $R_{out-long}$. 
            Это объясняется \"поворотом\" 2-х мерной проекции
            кореляционной функции в проекции \texttt{out-long}

            \begin{figure}[H]
                \centering
                \includegraphics[width=\linewidth]{images/all_out-long_2d_histos_centr_0-10_pos.pdf}
            \end{figure}
            \begin{figure}[H]
                \centering
                \includegraphics[width=\linewidth]{images/all_out-long_2d_histos_centr_10-30_pos.pdf}
                \caption{
                    "Поворот" проекций кореляционной функции для различных классов центральности $\pi^{+}$
                }
            \end{figure}

            \begin{figure}[H]
                \centering
                \includegraphics[width=\linewidth]{images/all_out-long_2d_histos_centr_0-10_neg.pdf}
            \end{figure}
            \begin{figure}[H]
                \centering
                \includegraphics[width=\linewidth]{images/all_out-long_2d_histos_centr_10-30_neg.pdf}
                \caption{\"Поворот\" проекций кореляционной функции для различных классов центральности $\pi^{-}$}
            \end{figure}

        \subsubsection{Отношение проекций кореляционных функций}
            Отсутствие выбросов или отклонений (в пределах статистической погрешности) 
            позволяет говорить об отсутствии зарядовой ассиметрии 

            \begin{figure}[H]
                \centering
                \includegraphics[width=\linewidth]{images/mr.png}
                \caption{
                    Отношение проекций для кореляционных функций пионов.
                }\label{fig:placeholder}
            \end{figure}

\section{Заключение}
    В рамках 9 семестра было проведено исследование пространственно-временных характеристик 
    источника пионов в столкновениях $^{197}\text{Au} + ^{197}\text{Au}$ при энергии $\sqrt{s_{NN}} = 3$~ГэВ 
    на основе данных транспортного моделирования \texttt{SMASH}.

    \textbf{Основные результаты работы:}
    \begin{enumerate}[leftmargin=*]
        \item \textbf{Разработан и верифицирован конвертер данных} 
            из формата \texttt{OSCAR2013} в формат \texttt{ROOT} (\texttt{McDst}). 
            Алгоритм обеспечивает корректное извлечение координат кинетического вымораживания 
            и идентификацию спектаторов. Валидация через сравнение распределений четырехвекторов 
            энергии-импульса показала полное совпадение с эталонными данными \texttt{SMASH} 
            (отношение гистограмм равно единице).

        \item \textbf{Построены одно- и двумерные корреляционные функции} 
        для $\pi^+\pi^+$ и $\pi^-\pi^-$ пар в интервале быстрот 
        $y \in [-1; 1]$ и для классов центральности $0\%$--$10\%$, $10\%$--$30\%$. 
        Аппроксимация функций гауссовой параметризацией позволила извлечь фемтоскопические параметры: 
        радиусы $R_{\text{out}}$, $R_{\text{side}}$, $R_{\text{long}}$, $R_{\text{out-long}}$ 
        и параметр хаотичности $\lambda$.

        \item \textbf{Исследованы зависимости радиусов и $\lambda$ от быстроты.} 
            Установлено, что:
            \begin{itemize}
                \item Фемтоскопические радиусы $R_{\text{out}}$ и $R_{\text{long}}$ 
                    минимальны в области центральных быстрот ($y \approx 0$) 
                    и квадратично возрастают с увеличением $|y|$. 
                    $R_{\text{side}}$ слабо зависит от парной быстроты;
                \item Параметр $\lambda$ демонстрирует аналогичное поведение, 
                    снижаясь в центральной области, что может указывать на вклад долгоживущих резонансов;
                \item Наблюдается рост кросс-термина $R_{\text{out-long}}$ 
                    с удалением от центральной быстроты, что интерпретируется как 
                    эффект «поворота» корреляционной функции в плоскости 
                    ($q_{\text{out}}$,$q_{\text{long}}$).
            \end{itemize}

        \item \textbf{Проанализировано отношение корреляционных функций} 
            для $\pi^+$ и $\pi^-$. В пределах текущей статистической точности 
            значимых отклонений от единицы не обнаружено, что свидетельствует 
            об отсутствии заметной зарядовой асимметрии в фемтоскопических радиусах 
            для системы Au+Au при $\sqrt{s_{NN}} = 3$~ГэВ в рамках модели \texttt{SMASH}.
    \end{enumerate}

    \newpage

\addcontentsline{toc}{chapter}{Литература}
\printbibliography[title = Литература]
\end{document}